% ----------------------------------------------------------------- thank you
% The title slide's layout: its panel on the right, the thanks where the title was; drawn inside the
% frame, so that the notes' thumbnail shows it too.
\begin{frame}[plain,noframenumbering]
\vbox to 0pt{\vskip-2mm\hbox to 0pt{\hskip-8mm\begin{tikzpicture}[x=1mm,y=1mm]
  \useasboundingbox (0,0) rectangle (160,-90);
  \titlepanel
  \node[anchor=north west,inner sep=0pt] at (8,-8)
    {\includegraphics[width=30mm,height=12mm,keepaspectratio]{\institutionlogo}};
  \node[anchor=north west,inner sep=0pt,text width=88mm,align=left] at (8,-36) {%
    {\fontsize{32}{36}\selectfont\bfseries\color{primary}Thank you\par}\vskip4mm
    \textcolor{accent}{\rule{14mm}{0.8pt}}\par\vskip3mm
    {\fontsize{11}{13}\selectfont\color{ink}\insertauthor\par}};
\end{tikzpicture}\hss}\vss}
\par\rule{0pt}{80mm}
\note{\textbf{8:10--8:30}\par
    Thank you.\par
    Questions to open the discussion: would the edge term help in our own work; how would we check it
      on several runs; and what would make us trust a sharper output.\par
    Backup slides follow: the loss in full and the references.}
\end{frame}
